\begin{tikzpicture}
\begin{scope}[ultra thick]
    \draw (0, 0) node[above] {proxy-client/server} -- ++(0, -7);
    \draw (6.5, 0) node[above,align=center] {SOCKS\\proxy} -- ++(0, -7);
    \draw (13, 0) node[above] {remote-client} -- ++(0, -7);
\end{scope}
\tikzset{
    msg/.style={font=\small,inner sep=0.2mm},
    send/.style={thick,-Latex},
    send misc/.style={thick,dotted,-Latex},
}
\draw[send] (0, -.5) -- ++(6.5, 0) node[midway,above,msg] {open TCP connection};
%\draw[send] (6.5, -.75) -- ++(6.5, 0);
\draw[red,send] (0, -1) -- ++(6.5, 0) node[midway,above,msg] {BIND 0:0 \textit{or} BIND IP:port};
\draw[red,send] (6.5, -1.5) -- ++(-6.5, 0) node[midway,above,msg] {reply w/ IP+port on proxy};
\draw[red,send] (13, -2) -- ++(-6.5, 0) node[midway,above,msg] {open TCP connection};
%\draw[send] (6.5, -2.25) -- ++(6.5, 0) node[midway,above,msg] {open TCP connection};
\draw[send] (6.5, -2.5) -- ++(-6.5, 0) node[midway,above,msg] {reply w/ success status};
%
\draw[ultra thick,red] (-.5, -4) rectangle (13.5, -7);
    \node[align=left,anchor=south,font=\small,red] at (4, -4) { now proxy-client/remote-client can use TCP connection \\ as if directly connected };
\draw[send] (0, -4.5) -- ++(6.5, 0) node[midway,above,msg] {data for server};
\draw[send] (6.5, -4.75) -- ++(6.5, 0) node[midway,above,msg] {data for server};
\draw[send] (13, -5.5) -- ++(-6.5, 0) node[midway,above,msg] {data from server};
\draw[send] (6.5, -5.75) -- ++(-6.5, 0) node[midway,above,msg] {data from server};
\end{tikzpicture}
